% appendix-a-simulation-dataset.tex -- part of thesis_v3.tex; do not add \documentclass here.
% Draft with the thesis-chapter workflow. Unsupported claims must use \fillin{...}.

\chapter{Power Integrity Simulation Environment and Dataset Specification}
\label{ch:appendix_simulation_dataset}

This appendix is reference material. It records the simulation environment that produced the
ground-truth labels, the file formats that carry them, and the on-disk specification of the
datasets used for training and for active learning. Tables are preferred to prose throughout;
each table is preceded by one short orienting paragraph. Everything stated here is taken from
repository artifacts (source code, configuration files, run reports, dataset metadata) as read on
2026-08-05, except Section~\ref{sec:appA_multitype}, whose figures are read from the Windows
simulation host --- a progress file, per-band output folders and a running process --- and carry
the observation timestamps given with them. Quantities that exist only inside the eCADSTAR project
on the Windows host --- solver settings, stack-up, port geometry, board dimensions, capacitor
models for the binary dataset --- are marked as gaps rather than reconstructed.

\section{Simulation environment}
\label{sec:appA_sim_env}

\subsection{Solver and invocation}
\label{subsec:appA_solver}

All power-integrity labels are produced by the eCADSTAR PI/EMI analysis engine, driven in batch
mode from WSL. There is no GUI automation path: the AutoHotkey driver was removed and the headless
CLI is the only route to the engine, with no fallback and no toggle to switch back
(\texttt{docs/ecadstar\_headless\_cli.md}; \texttt{active\_learning\_pi/al/ecadstar.py}). One
residual key remains in the primary active-learning configuration, a null \texttt{ahk\_exe} entry
under \texttt{ecadstar}, which is no longer read by any driver.

\begin{table}[htbp]
\centering
\caption{Solver executable and command-line interface.}
\label{tab:appA_solver}
\begin{tabular}{p{0.30\textwidth} p{0.60\textwidth}}
\hline
Item & Value \\
\hline
Engine executable &
\texttt{C:\textbackslash{}Program Files\textbackslash{}eCADSTAR\textbackslash{}eCADSTAR
2023.0\textbackslash{}Analysis\textbackslash{}bin\textbackslash{}engineer.exe} \\
Argument line &
\texttt{"<erf>" -{}-batch "<peb>" -{}-batch-auto-exit}, optionally extended with
\texttt{-{}-impulse-port <n>} \\
\texttt{-b} / \texttt{-{}-batch <file>} & execute the given batch file after the design loads \\
\texttt{-be} / \texttt{-{}-batch-auto-exit} & quit the application when the batch finishes \\
\texttt{-p} / \texttt{-{}-impulse-port <n>} & IMPULSE model-library server port \\
Positional argument & the design riffile (\texttt{.erf} / \texttt{.rif}) \\
\hline
\end{tabular}
\end{table}

The launcher (\texttt{pipelines/dataset\_sim/ecadstar.py}, function
\texttt{run\_ecadstar\_batch\_headless}) starts the engine through PowerShell
\texttt{Start-Process -PassThru -NoNewWindow} and then blocks on \texttt{WaitForExit(timeout\_ms)}.
Return codes are \texttt{0} for a clean self-exit, \texttt{124} when the timeout expires and the
process is killed, and \texttt{-1} when PowerShell is unavailable. The default timeout is
\texttt{172\_800} seconds. Windows paths held in the configuration files are converted from WSL
form by \texttt{resolve\_windows\_path} / \texttt{windows\_path\_str}.

\subsection{Design under analysis and host directories}
\label{subsec:appA_design_paths}

A single board design is simulated throughout. All of the paths below are absolute Windows paths;
the simulation half of this work is therefore reproducible only on that host.

\begin{table}[htbp]
\centering
\caption{Design, power bus and host directories used by every simulation pipeline.}
\label{tab:appA_paths}
\begin{tabular}{p{0.28\textwidth} p{0.62\textwidth}}
\hline
Item & Value \\
\hline
Design file &
\texttt{C:\textbackslash{}Users\textbackslash{}muthusamy\textbackslash{}Desktop\textbackslash{}design\textbackslash{}H-shape.emc\textbackslash{}H-shape.erf} \\
EMC output directory &
\texttt{C:\textbackslash{}Users\textbackslash{}muthusamy\textbackslash{}Desktop\textbackslash{}design\textbackslash{}H-shape.emc}
(engine writes \texttt{PI-1} \ldots{} \texttt{PI-N}) \\
PEB staging directory &
\texttt{C:\textbackslash{}Users\textbackslash{}muthusamy\textbackslash{}Desktop\textbackslash{}design\textbackslash{}PEB} \\
Raw output root &
\texttt{C:\textbackslash{}Users\textbackslash{}muthusamy\textbackslash{}Desktop\textbackslash{}Raw}
(\texttt{RAW\_ROOT\_WIN}) \\
PEB build directory &
\texttt{C:\textbackslash{}Users\textbackslash{}muthusamy\textbackslash{}Desktop\textbackslash{}Raw\textbackslash{}peb}
(\texttt{PEB\_DIR\_WIN}) \\
Analysed power bus & \texttt{Power\_GND} \\
Spectrum export port(s) & \texttt{IC1\_Port1,IC2\_Port2} (combinations pipeline);
\texttt{IC1\_Port1} (multi-type pipeline and the AL PEB builder) \\
\hline
\end{tabular}
\end{table}

The two port strings are a configuration difference between pipelines, not a defect found in an
artifact. \fillin{state which port convention applies to the 231-bin spectra actually stored in
\texttt{datasets/data\_multifreq\_train}; settle it by inspecting the harvested CSV filenames for a
legacy PI folder against \texttt{peb\_components} in the pipeline that produced it}

\subsection{Solver settings not recorded in the repository}
\label{subsec:appA_solver_gaps}

The following belong to the eCADSTAR project itself and were not found in any tracked file. They
are listed explicitly so that the appendix does not imply knowledge it does not have.

\begin{table}[htbp]
\centering
\caption{Simulation parameters that live in the eCADSTAR project on the Windows host and are not
recoverable from the repository.}
\label{tab:appA_solver_gaps}
\begin{tabular}{p{0.38\textwidth} p{0.52\textwidth}}
\hline
Parameter group & Status \\
\hline
Frequency sweep definition, number of sweep points &
\fillin{export the PI analysis setup from \texttt{H-shape.erf} and record the sweep definition; the
only repository trace is the 231-point axis in \texttt{configs/Frequency\_data\_hz.npy} and the
\texttt{EXPECTED\_IMP\_LENGTH = 231} check} \\
Mesh / plane-pair settings &
\fillin{record from the eCADSTAR PI analysis setup dialog} \\
Port definitions and locations (\texttt{IC1\_Port1}, \texttt{IC2\_Port2}) &
\fillin{record port placement and reference from the design} \\
Layer stack-up, dielectric and copper material parameters &
\fillin{export stack-up table from \texttt{H-shape.erf}} \\
Board outline dimensions &
\fillin{record board size; needed to give the $64\times64$ heatmap grid a physical scale} \\
Decap C / ESR / ESL for the single-type (binary) dataset &
\fillin{record the capacitor model attached to slots \texttt{C1}--\texttt{C52} in the binary
campaign; the \texttt{TYPE\_CATALOG} in
\texttt{scrap/generation/generate\_peb\_multitype.py} covers only the multi-type pipeline} \\
\hline
\end{tabular}
\end{table}

\figplace{Annotated board view of H-shape.erf showing the 52 decap slots, the IC port(s) and the power-bus outline.}

\section{Batch file format and simulation phases}
\label{sec:appA_peb}

\subsection{PEB structure}
\label{subsec:appA_peb_structure}

A PEB file is XML. Its root element is
\texttt{<Batch Analysis="PI" Version="2" PowerBus="\{powerbus\}">}, and it contains an
\texttt{<Iterate>} list of \texttt{<Group>} blocks, one group per layout row. Each group first
writes the 52 component enable lines, then the analysis command for that group.

\begin{table}[htbp]
\centering
\caption{Elements emitted inside a PEB group.}
\label{tab:appA_peb_elements}
\begin{tabular}{p{0.26\textwidth} p{0.64\textwidth}}
\hline
Purpose & Emitted line \\
\hline
Slot occupancy (52 lines) &
\texttt{<Edit Table="Component" Name="C\{i+1\}" Column="Name" Value="true|false"/>} for
\texttt{C1}\ldots\texttt{C52} \\
PI-Distribution group &
\texttt{<Loop Analysis="PI-Distribution" Count="1"><EditPIDistribution
Frequency="\{freq\}"/></Loop>} \\
PI-Spectrum group &
\texttt{<CreatePISpectrum Component="\{components\}" Format="ARV,CSV" />} \\
Frequency string format &
\texttt{f"\{int(round(float(mhz)))\}e6"} (Hz) \\
\hline
\end{tabular}
\end{table}

For the multi-type campaign the per-component block grows from one line to four: Component/Name,
Component/Value (capacitance), PI Decap/ESR and PI Decap/ESL. Empty slots write \texttt{false} and
a value of \texttt{0}.

\subsection{Two-phase dataset simulation}
\label{subsec:appA_two_phase}

Dataset simulation is split into an impedance phase and a per-band distribution phase. Both phases
iterate over the same layout rows, and each produces one \texttt{PI-$n$} output folder per row.

\begin{table}[htbp]
\centering
\caption{Simulation phases and their harvested artifacts.}
\label{tab:appA_phases}
\begin{tabular}{p{0.20\textwidth} p{0.34\textwidth} p{0.36\textwidth}}
\hline
Phase & PEB builder & Harvested file \\
\hline
Impedance &
\texttt{build\_impedance\_peb}, spectrum only (\texttt{include\_distribution=False}), one PEB over
all rows &
first \texttt{*.csv} in the PI folder whose name contains \texttt{ic1} (case-insensitive) \\
Distribution &
\texttt{build\_distribution\_peb}, distribution only (\texttt{include\_spectrum=False}), one PEB
per MHz band &
\texttt{<PI-n>/Power\_GND/Z\_\{MHz:04d\}.000MHz.map} \\
\hline
\end{tabular}
\end{table}

After a batch, \texttt{wait\_for\_pi\_outputs} polls the EMC directory until \texttt{PI-1} through
\texttt{PI-N} are present, checking for \texttt{.map} files when the mode is \texttt{distribution}
and for \texttt{.csv} files when the mode is \texttt{impedance}. CSV row $i$ of the layout list maps
to \texttt{PI-$(i+1)$} in eCADSTAR script order. In the merged combination list the legacy rows keep
the low indices, \texttt{PI-1}\ldots\texttt{PI-19499}, and the newer combinations follow at
\texttt{PI-19500}\ldots\texttt{PI-29499}.

\subsection{Pipeline configuration constants}
\label{subsec:appA_pipeline_config}

The runnable pipelines take no command-line arguments; their behaviour is fixed by a configuration
block of module-level constants. The values below are those present in the block at the time of
reading.
\fillin{establish whether \texttt{SWEEP\_MHZ} was edited between campaigns and which band sets were
actually simulated; \texttt{logs/append\_queue/done.jsonl} records the bands that completed}

\begin{table}[htbp]
\centering
\caption{Single-type combinations simulation pipeline
(\texttt{pipelines/dataset\_sim/run\_combinations\_sim\_pipeline.py}).}
\label{tab:appA_comb_pipeline}
\begin{tabular}{p{0.40\textwidth} p{0.50\textwidth}}
\hline
Constant & Value \\
\hline
\texttt{SWEEP\_MHZ} & \texttt{[170, 180, 280, 430, 470, 490, 550, 590]} \\
Input CSV & \texttt{data/heatmaps/peb\_with\_29k/all\_combinations\_merged.csv} \\
\texttt{ECADSTAR\_WAIT\_TIMEOUT\_SEC} & \texttt{172\_800} \\
\texttt{ECADSTAR\_WAIT\_POLL\_SEC} & \texttt{120} \\
\hline
\end{tabular}
\end{table}

\fillin{note whether the 490, 550 and 590\,MHz bands in \texttt{SWEEP\_MHZ} ever reached the
dataset --- they are absent from the 24 anchors in
\texttt{datasets/data\_multifreq\_train/dataset\_meta.json}; the append queue in
\texttt{logs/append\_queue/done.jsonl} would settle it}

\subsection{Active-learning batches}
\label{subsec:appA_al_batches}

The active-learning loop exports all selected candidates of a cycle into one combined PEB
(\texttt{active\_learning\_pi/al/peb\_batch.py}, \texttt{build\_peb\_from\_selection}), carrying a
per-sample frequency formatted as \texttt{"\{int(round(mhz))\}e6"}. The builder sets
\texttt{include\_distribution=True} and \texttt{include\_spectrum = not heatmap\_only}. In the
primary configuration both \texttt{heatmap\_only\_peb} and \texttt{heatmap\_only\_labels} are
\texttt{true}, so AL batches request PI-Distribution output only.

\begin{table}[htbp]
\centering
\caption{eCADSTAR settings in the primary AL configuration
(\texttt{active\_learning\_pi/config/exp059\_gp\_error.json}, key \texttt{ecadstar}).}
\label{tab:appA_al_ecadstar}
\begin{tabular}{p{0.45\textwidth} p{0.30\textwidth}}
\hline
Key & Value \\
\hline
\texttt{clear\_lock\_file} & \texttt{true} \\
\texttt{wait\_for\_pi} & \texttt{true} \\
\texttt{wait\_timeout\_sec} & \texttt{18000} \\
\texttt{wait\_batch\_start\_timeout\_sec} & \texttt{900} \\
\texttt{wait\_poll\_sec} & \texttt{40} \\
\hline
\end{tabular}
\end{table}

Label ingestion (\texttt{active\_learning\_pi/al/ingest\_labels.py}) reads \texttt{PI-1} through
\texttt{PI-N} from the EMC directory, copies the \texttt{.map} file for each candidate's frequency
into \texttt{labels\_dir/sample\_\{i:03d\}/}, and writes \texttt{occupancy.npy},
\texttt{pi\_freq\_mhz.npy} and \texttt{heatmap.npy}, the last produced by
\texttt{create\_Heatmaps} at grid size 64.

\section{From solver output to tensors}
\label{sec:appA_encoding}

\subsection{Decap slots and occupancy vectors}
\label{subsec:appA_occupancy}

Decap slots are arranged on a $7\times8$ board grid. Four cells are declared invalid ---
$(0,3)$, $(0,4)$, $(6,3)$ and $(6,4)$ --- leaving $7\times8-4 = 52$ valid positions, labelled
\texttt{C1} through \texttt{C52}. A placement is therefore the binary vector
$\mathbf{b}\in\{0,1\}^{52}$, with budget $K = \sum_i b_i$. On disk an occupancy vector is a
length-52 \texttt{float32} array; it is binarised with a threshold of $0.5$.

\figplace{The 7x8 slot grid with the four invalid cells shaded and the C1--C52 labelling order.}

\subsection{Impedance spectra}
\label{subsec:appA_impedance}

The exported PI-Spectrum CSV is read with delimiter \texttt{;}, \texttt{skiprows=4} and
\texttt{usecols=[1]}. A vector is rejected unless its length is exactly 231
(\texttt{EXPECTED\_IMP\_LENGTH = 231}). The impedance $Z(f)$ used throughout the thesis is this
231-bin vector.

\begin{table}[htbp]
\centering
\caption{Frequency axis and target curve stored under \texttt{configs/}.}
\label{tab:appA_freq_axis}
\begin{tabular}{p{0.34\textwidth} p{0.56\textwidth}}
\hline
Artifact & Contents \\
\hline
\texttt{configs/Frequency\_data\_hz.npy} &
shape \texttt{(231,)}, minimum \texttt{1000000.0}\,Hz, maximum \texttt{600000000.0}\,Hz;
first values \texttt{1000000, 1100000, 1200000, 1300000, 1400000}; last values
\texttt{5.8e+08, 5.9e+08, 6.0e+08} \\
\texttt{configs/target\_impedance.npy} &
shape \texttt{(231, 1)}, minimum \texttt{0.8}, maximum \texttt{49.99999999999999} \\
\hline
\end{tabular}
\end{table}

Neither file carries metadata about its own construction.
\fillin{document how the 231-point frequency grid was generated (the \texttt{.npy} has no
provenance); the eCADSTAR sweep definition in \texttt{H-shape.erf} would settle it}
\fillin{document how $Z_{\mathrm{target}}(f)$ in \texttt{configs/target\_impedance.npy} was
specified --- per-band values, derivation, and the design rule behind the 0.8 and 50 limits}

\subsection{PI-distribution heatmaps}
\label{subsec:appA_heatmap}

Each \texttt{.map} file is parsed and interpolated onto a fixed $64\times64$ grid. The steps
implemented in \texttt{libs/data\_creation/heatmap.py} are listed below in the order applied.

\begin{table}[htbp]
\centering
\caption{Heatmap construction from a PI-Distribution \texttt{.map} file.}
\label{tab:appA_heatmap_steps}
\begin{tabular}{p{0.34\textwidth} p{0.56\textwidth}}
\hline
Step & Setting \\
\hline
Coordinate handling & shifted to the minimum, scaled by \texttt{1e-5} \\
Value floor & $z$ values floored at \texttt{0.01} \\
Duplicate vertices & averaged \\
Interpolation & \texttt{RBFInterpolator}, kernel \texttt{thin\_plate\_spline},
\texttt{smoothing=0} \\
Fallback & cubic \texttt{griddata} with nearest-neighbour NaN fill \\
Grid size & \texttt{64} \\
Stored tensor & \texttt{np.stack([Zi * mask\_board, mask\_board], axis=0)}: channel 0 is the masked
interpolated impedance, channel 1 is the mask \\
Board mask & \texttt{configs/binary\_mask.npy}, \texttt{(64, 64)} boolean, \texttt{3432} True
pixels \\
\hline
\end{tabular}
\end{table}

\section{Dataset layout and metadata}
\label{sec:appA_dataset_layout}

\subsection{Directory structure}
\label{subsec:appA_dir_structure}

Datasets use a layout-store organisation: per-layout quantities are written once under a design
directory, and per-sample quantities (one per layout--frequency pair) are written under flat
sample directories.

\begin{table}[htbp]
\centering
\caption{Layout-store dataset structure.}
\label{tab:appA_store}
\begin{tabular}{p{0.34\textwidth} p{0.56\textwidth}}
\hline
Path & Contents \\
\hline
\texttt{layouts/\{design\_id\}/occ.npy} & occupancy vector, shape \texttt{(52,)} \\
\texttt{layouts/\{design\_id\}/imp.npy} & impedance, reshaped to \texttt{(-1, 1)}, i.e.
\texttt{(231, 1)} in the raw store \\
\texttt{heatmap/sample\_N.npy} & the 2-channel $64\times64$ stack \\
\texttt{PI\_freq/sample\_N.npy} & the sample's frequency in Hz, \texttt{float64} \\
\texttt{manifest.csv} & one row per sample \\
\texttt{dataset\_meta.json} & dataset-level metadata \\
\hline
\end{tabular}
\end{table}

The manifest columns are \texttt{sample\_name}, \texttt{design\_id}, \texttt{freq\_label},
\texttt{freq\_mhz}, \texttt{freq\_hz}, \texttt{source\_folder}, \texttt{append\_tag},
\texttt{pi\_number} and \texttt{decap\_index} --- nine columns in total.

\subsection{\texttt{dataset\_meta.json} fields}
\label{subsec:appA_meta_fields}

\begin{table}[htbp]
\centering
\caption{Fields written by \texttt{libs/dataset\_meta.py}.}
\label{tab:appA_meta_fields}
\begin{tabular}{p{0.30\textwidth} p{0.60\textwidth}}
\hline
Field & Notes \\
\hline
\texttt{schema\_version} & integer \\
\texttt{stage} & \texttt{raw} or \texttt{normalized} \\
\texttt{source\_script} & the script that produced the dataset \\
\texttt{generated\_at\_utc} & ISO-8601 UTC timestamp \\
\texttt{dataset\_root} & absolute path to the dataset root, as written by
\texttt{libs/dataset\_meta.py} line 172 \\
\texttt{storage} & \texttt{layout\_store} \\
\texttt{counts} & \texttt{manifest\_rows}, \texttt{unique\_layouts}, \texttt{heatmap\_files},
\texttt{pi\_freq\_files}, \texttt{layout\_directories} \\
\texttt{size} & \texttt{total\_bytes}, \texttt{total\_mb}, \texttt{by\_subdirectory\_mb} \\
\texttt{pi\_frequencies\_mhz} & the anchor list \\
\texttt{samples\_per\_mhz}, \texttt{samples\_per\_freq\_label} & per-band coverage \\
\hline
\multicolumn{2}{l}{\emph{Additional fields on normalized sets:}} \\
\texttt{normalization} & \texttt{stats\_file}, \texttt{has\_heatmap\_stats} \\
\texttt{mode} & e.g.\ \texttt{full} \\
\texttt{source\_dataset} & the raw set it derives from; also written as an absolute path \\
\texttt{heatmap\_norm} & normalization mode name \\
\texttt{max\_k\_filter} & budget cut-off applied \\
\hline
\end{tabular}
\end{table}

\subsection{Dataset inventory and version control}
\label{subsec:appA_inventory}

\textbf{None of the dataset directories are under version control.} The repository
\texttt{.gitignore} ignores \texttt{datasets/}, \texttt{data\_multi\_norm\_robust/},
\texttt{checkpoints/}, \texttt{**/checkpoints/}, \texttt{*.pt} and \texttt{*.pth};
\texttt{git check-ignore} confirms that
\texttt{datasets/data\_multifreq\_train/manifest.csv} and
\texttt{datasets/data\_multifreq\_train\_norm\_unbounded/normalization\_stats.json} are matched by
\texttt{.gitignore} line 20 (\texttt{datasets/}). Reproducing any result in this thesis therefore
requires rebuilding the datasets from the raw eCADSTAR output on the Windows host, or obtaining the
directories separately.

\begin{table}[htbp]
\centering
\caption{Dataset directories present on the working machine (all git-ignored). Paths are absolute.}
\label{tab:appA_inventory}
\begin{tabular}{p{0.56\textwidth} p{0.34\textwidth}}
\hline
Path & Role \\
\hline
\texttt{/home/ubuntu/genai\_pdn/datasets/data\_multifreq\_train} & raw training set \\
\texttt{/home/ubuntu/genai\_pdn/datasets/data\_multifreq\_train\_norm\_unbounded} & normalized set;
training input for exp059 \\
\texttt{/home/ubuntu/genai\_pdn/datasets/data\_multifreq\_train\_norm\_robust} & bounded robust
normalization, earlier experiments \\
\texttt{/home/ubuntu/genai\_pdn/datasets/data\_multifreq\_train\_K5} & $K\le5$ subset \\
\texttt{/home/ubuntu/genai\_pdn/datasets/data\_multifreq\_al\_overlay\_exp057} & AL overlay
(legacy) \\
\texttt{/home/ubuntu/genai\_pdn/datasets/data\_multifreq\_al\_overlay\_exp058} & AL overlay
(legacy) \\
\texttt{/home/ubuntu/genai\_pdn/datasets/data\_multifreq\_al\_overlay\_exp059} & AL overlay
(current track) \\
\texttt{/home/ubuntu/genai\_pdn/data\_multi\_norm\_robust} & older backup at the repository root \\
\hline
\end{tabular}
\end{table}

\texttt{datasets/data\_multifreq\_al\_overlay\_exp060} is referenced by
\texttt{experiments/exp060\_multitype\_occ/config.yaml} line 23 but does not exist on disk.

The current model track, \texttt{exp059\_capacity\_freq}, trains from \texttt{data\_dir}
\texttt{datasets/data\_multifreq\_train\_norm\_unbounded}. The primary AL configuration uses the
same \texttt{data\_dir}, \texttt{normalization\_stats\_json}
\texttt{datasets/data\_multifreq\_train\_norm\_unbounded/normalization\_stats.json} and
\texttt{overlay\_data\_dir} \texttt{datasets/data\_multifreq\_al\_overlay\_exp059}.

\fillin{state the train/validation/test split definition for the normalized dataset --- no split
artifact was located in the dataset directories read for this appendix; the split logic in the
exp059 dataloader would settle it}

\section{Dataset statistics}
\label{sec:appA_dataset_stats}

\subsection{Layout pool}
\label{subsec:appA_layout_pool}

The layout pool is a merged combination list. The merge report
(\texttt{data/heatmaps/merged\_combinations\_report.json}, generated
\texttt{2026-07-01T13:13:39.245630+00:00}) records \texttt{n\_old} 19499 rows from
\texttt{all\_combinations.csv}, \texttt{n\_new} 10000 rows from \texttt{combinations.csv},
\texttt{overlap\_rows} 0, giving 29499 merged rows. The merged CSV read by the simulation pipeline
is \texttt{data/heatmaps/peb\_with\_29k/all\_combinations\_merged.csv}, which has 29499 lines.
The output path recorded in the merge report itself
(\texttt{data/heatmaps/all\_combinations\_merged.csv}) no longer exists at that location; see
Section~\ref{sec:appA_inconsistencies}.

\subsection{Raw training set}
\label{subsec:appA_raw_set}

This dataset root is written to by the append worker, which was still being triggered while this
appendix was assembled; the values below are the state at the read time given in the table and must
be re-read before submission. The counts have not changed across the re-reads made for this
appendix, but \texttt{generated\_at\_utc} has.

\begin{table}[htbp]
\centering
\caption{\texttt{datasets/data\_multifreq\_train} (raw), as read on 2026-08-06.}
\label{tab:appA_raw}
\begin{tabular}{p{0.40\textwidth} p{0.50\textwidth}}
\hline
Field & Value \\
\hline
\texttt{schema\_version} & 1 \\
\texttt{stage} & \texttt{raw} \\
\texttt{source\_script} &
\texttt{pipelines/data/append\_merged\_combinations\_multifreq.py} \\
\texttt{generated\_at\_utc} & \texttt{2026-08-05T19:28:42.489387+00:00} \\
\texttt{dataset\_root} & \texttt{/home/ubuntu/genai\_pdn/\allowbreak datasets/data\_multifreq\_train} \\
\texttt{storage} & \texttt{layout\_store} \\
\texttt{manifest\_rows} & 688477 \\
\texttt{unique\_layouts} & 29499 \\
\texttt{heatmap\_files} & 688477 \\
\texttt{pi\_freq\_files} & 688477 \\
\texttt{layout\_directories} & 29499 \\
\texttt{total\_mb} & 43307.27 \\
\quad heatmap & 43113.86 \\
\quad layouts & 39.05 \\
\quad PI\_freq & 89.3 \\
\texttt{samples\_per\_mhz} & 29499 at every anchor except 80\,MHz, which has 10000 \\
\hline
\end{tabular}
\end{table}

\fillin{explain why the 80\,MHz band covers only 10000 of the 29499 layouts --- no artifact was
located that records the reason; the per-band append log in \texttt{logs/append\_queue/done.jsonl}
and the 80\,MHz raw folder count would settle it}

\subsection{Frequency anchors}
\label{subsec:appA_anchors}

The anchor list is the set of inspection frequencies at which PI-Distribution heatmaps exist.
\texttt{datasets/data\_multifreq\_train/dataset\_meta.json} lists 24 anchors, in MHz:
10.0, 63.0, 80.0, 100.0, 120.0, 150.0, 170.0, 180.0, 200.0, 230.0, 250.0, 270.0, 280.0, 300.0,
330.0, 350.0, 370.0, 390.0, 400.0, 420.0, 430.0, 450.0, 470.0, 500.0.

Anchor resolution at load time follows a four-step order, per \texttt{docs/dataset.md}:
\texttt{dataset\_meta.json} \texttt{pi\_frequencies\_mhz}, then \texttt{manifest.csv}
\texttt{freq\_mhz}, then \texttt{datasets/data\_multifreq\_train/dataset\_meta.json}, then the
legacy YAML defaults. That documentation statement was not verified against the implementation in
\texttt{src\_vae/others/multifreq\_anchors.py} line by line.

Three sources give three different anchor counts. \texttt{docs/dataset.md} lines 96--101 describe a
``Full 16-MHz set'' and ``Same 16 MHz in combinations sim batches''; the legacy fallback
\texttt{configs/multifreq\_anchors.yaml} lists 20 anchors (10, 63, 80, 100, 120, 150, 170, 200,
230, 250, 270, 300, 330, 350, 370, 390, 400, 420, 450, 500), omitting 180, 280, 430 and 470\,MHz;
\texttt{dataset\_meta.json} lists the 24 above. Per the resolution order in
\texttt{docs/dataset.md} lines 42--51, the \texttt{dataset\_meta.json} list is what is actually
loaded, and \texttt{configs/multifreq\_anchors.yaml} is declared a fallback used only when no
metadata or manifest is present.
\fillin{resolve the conflict between the ``16-MHz set'' description in \texttt{docs/dataset.md}
lines 96--101, the 20-anchor list in \texttt{configs/multifreq\_anchors.yaml}, and the 24-anchor
list in \texttt{datasets/data\_multifreq\_train/dataset\_meta.json}; decide which is stale and
correct the documentation}

\fillin{document the provenance of the 24-anchor selection --- no artifact was located that records
why these particular frequencies were chosen}

\subsection{Normalized training set}
\label{subsec:appA_norm_set}

\begin{table}[htbp]
\centering
\caption{\texttt{datasets/data\_multifreq\_train\_norm\_unbounded} (normalized; the exp059 training
input).}
\label{tab:appA_norm}
\begin{tabular}{p{0.40\textwidth} p{0.50\textwidth}}
\hline
Field & Value \\
\hline
\texttt{stage} & \texttt{normalized} \\
\texttt{source\_script} & \texttt{pipelines/normalize/multifreq.py} \\
\texttt{generated\_at\_utc} & \texttt{2026-07-06T13:39:03.768907+00:00} \\
\texttt{mode} & \texttt{full} \\
\texttt{source\_dataset} & \texttt{/home/ubuntu/genai\_pdn/\allowbreak datasets/data\_multifreq\_train} \\
\texttt{heatmap\_norm} & \texttt{robust\_log1p\_per\_mhz\_unbounded} \\
\texttt{max\_k\_filter} & 30 \\
\texttt{manifest\_rows} & 582997 \\
\texttt{unique\_layouts} & 24979 \\
\texttt{heatmap\_files} & 582997 \\
\texttt{pi\_freq\_files} & 582997 \\
\texttt{layout\_directories} & 24979 \\
\texttt{total\_mb} & 9344.12 \\
\quad heatmap & 9180.49 \\
\quad layouts & 33.06 \\
\quad PI\_freq & 75.61 \\
\texttt{samples\_per\_mhz} & 24979 at every anchor except 80\,MHz, which has 8480 \\
\texttt{pi\_frequencies\_mhz} & the same 24 anchors as the raw set \\
\hline
\end{tabular}
\end{table}

The budget filter is applied by \texttt{pipelines/normalize/build\_train\_norm\_unbounded.py},
which sets \texttt{DATA\_DIR} to \texttt{datasets/data\_multifreq\_train}, \texttt{OUTPUT\_DIR} to
\texttt{datasets/data\_multifreq\_train\_norm\_unbounded} and \texttt{MAX\_K} to 30, so only layouts
with $K \le 30$ enter the normalized set. Note that this normalized set was generated on
2026-07-06, whereas the raw set it names as its source was last refreshed on 2026-08-05; the two
artifacts are not snapshots of the same moment, and the $K\le30$ filter accounts for only part of
the difference in counts.

\subsection{$K\le5$ subset}
\label{subsec:appA_k5}

\begin{table}[htbp]
\centering
\caption{\texttt{datasets/data\_multifreq\_train\_K5} and its dry-run report.}
\label{tab:appA_k5}
\begin{tabular}{p{0.40\textwidth} p{0.50\textwidth}}
\hline
Field & Value \\
\hline
\texttt{source\_script} & \texttt{pipelines/data/copy\_multifreq\_k\_subset.py} \\
\texttt{generated\_at\_utc} & \texttt{2026-08-02T17:03:40.514745+00:00} \\
\texttt{manifest\_rows} & 131766 \\
\texttt{unique\_layouts} & 5650 \\
\texttt{total\_mb} & 8288.11 \\
\texttt{pi\_frequencies\_mhz} & the same 24 anchors \\
\texttt{samples\_per\_mhz} & 5650 at every anchor except 80\,MHz, which has 1816 \\
\hline
\multicolumn{2}{l}{\emph{Dry-run report
(\texttt{datasets/data\_multifreq\_train\_K5\_dry\_run\_report.json},
\texttt{2026-08-02T17:01:03.238318+00:00}):}} \\
\texttt{max\_k} & 5 \\
\texttt{designs\_total} / \texttt{designs\_kept} & 29499 / 5650 \\
\texttt{manifest\_rows\_total} / \texttt{\dots\_kept} & 688477 / 131766 \\
\hline
\end{tabular}
\end{table}

The same report carries \texttt{k\_hist\_all}, the budget histogram of the full raw pool. Selected
entries: $K=2$: 1326; $K=3$: 1483; $K=4$: 1440; $K=5$: 1401; $K=20$: 912; $K=21$: 319; and 226 each
for $K=26$ through $K=50$.
\fillin{reproduce the complete \texttt{k\_hist\_all} histogram as a table or bar chart from
\texttt{datasets/data\_multifreq\_train\_K5\_dry\_run\_report.json} if the budget distribution is
discussed in the main text}

\subsection{Active-learning overlay (exp059)}
\label{subsec:appA_overlay}

\begin{table}[htbp]
\centering
\caption{\texttt{datasets/data\_multifreq\_al\_overlay\_exp059}, read from disk on 2026-08-05.}
\label{tab:appA_overlay}
\begin{tabular}{p{0.40\textwidth} p{0.50\textwidth}}
\hline
Item & Value \\
\hline
\texttt{manifest.csv} & 2001 lines (header plus 2000 rows) \\
Layout directories & 2000 \\
Subdirectories & \texttt{PI\_freq/}, \texttt{heatmap/}, \texttt{layouts/} \\
Other files & \texttt{normalization\_stats.json},
\texttt{al\_heatmap\_only\_design\_ids.json} \\
\texttt{dataset\_meta.json} & absent \\
\texttt{design\_id} form &
e.g.\ \texttt{al\_al\_exp059\_gp\_error\_001\_i0001\_c009} \\
\texttt{source\_folder} & \texttt{active\_learning} \\
Overlay \texttt{normalization\_stats.json} &
same key structure as the training set; \texttt{background\_value} \texttt{-3.1792},
\texttt{storage} \texttt{layout\_store}, \texttt{max\_k\_filter} 30 \\
\hline
\end{tabular}
\end{table}

Because the overlay has no \texttt{dataset\_meta.json}, its counts were read from the directory
listing rather than from a written metadata artifact.

\section{Normalization statistics}
\label{sec:appA_normalization}

\subsection{File structure}
\label{subsec:appA_norm_structure}

\texttt{normalization\_stats.json} has seven top-level keys: \texttt{background\_value},
\texttt{Heatmap}, \texttt{Impedance}, \texttt{PI\_freq}, \texttt{multifreq\_manifest},
\texttt{storage} and \texttt{max\_k\_filter}. It is assembled in
\texttt{pipelines/normalize/multifreq.py}.

\subsection{Heatmap normalization}
\label{subsec:appA_heatmap_norm}

Per-MHz heatmap statistics are computed on $\log(1+x)$ of the foreground pixels of that band. The
median and the interquartile range (IQR, $q_{75}-q_{25}$, floored at
\texttt{ROBUST\_IQR\_EPS = 1e-4}) define a robust $z$-score
$z = (\log(1+x_{\mathrm{fg}}) - \mathrm{median})/\mathrm{IQR}$. \texttt{clip\_min} is the 0.5th
percentile of $z$ and \texttt{clip\_max} the 99.5th; the background fill value is
$\mathrm{round}(\texttt{clip\_min} - 1.5, 4)$. In \emph{unbounded} mode the $z$-score is applied
without clipping to $[\texttt{clip\_min}, \texttt{clip\_max}]$; the clip is applied only when the
mode is not unbounded. The stored statistics still record the clip bounds.

\begin{table}[htbp]
\centering
\caption{Global heatmap statistics,
\texttt{datasets/data\_multifreq\_train\_norm\_unbounded/normalization\_stats.json}, key
\texttt{Heatmap}.}
\label{tab:appA_heatmap_stats}
\begin{tabular}{p{0.36\textwidth} p{0.50\textwidth}}
\hline
Key & Value \\
\hline
\texttt{norm\_mode} & \texttt{robust\_log1p\_per\_mhz\_unbounded} \\
\texttt{unbounded} & \texttt{true} \\
\texttt{clip\_min} & \texttt{-1.6791870594024658} \\
\texttt{clip\_max} & \texttt{6.275787830352783} \\
\texttt{z\_max} & \texttt{11.23223876953125} \\
\texttt{background\_value} & \texttt{-3.1792} \\
\texttt{log\_mean} & \texttt{0.0} \\
\texttt{log\_std} & \texttt{1.0} \\
\texttt{fg\_pixel\_count} & \texttt{2000845704} \\
\texttt{anchors\_mhz} & length 24 \\
\texttt{by\_mhz} & 24 keys \\
\hline
\end{tabular}
\end{table}

Each \texttt{by\_mhz} entry carries \texttt{median}, \texttt{iqr}, \texttt{clip\_min},
\texttt{clip\_max}, \texttt{background\_value}, \texttt{fg\_pixel\_count}, \texttt{z\_min} and
\texttt{z\_max}. The first four bands are reproduced below; the remaining 20 are in the same file
under the same key.

\begin{table}[htbp]
\centering
\caption{Per-MHz heatmap statistics, first four anchors.}
\label{tab:appA_heatmap_by_mhz}
\begin{tabular}{r r r r r}
\hline
MHz & median & iqr & clip\_min & clip\_max \\
\hline
10.0 & 0.05887173116207123 & 0.03866304084658623 & -0.7351016402244568 & 3.4719114303588867 \\
63.0 & 0.41589492559432983 & 0.23903468251228333 & -0.6570155620574951 & 3.969606876373291 \\
80.0 & 0.5136411190032959 & 0.29707127809524536 & -0.6531781554222107 & 6.275787830352783 \\
100.0 & 0.6201565265655518 & 0.357897512614727 & -1.2004997730255127 & 5.350724220275879 \\
\hline
\end{tabular}
\end{table}

\fillin{if the full per-band table is wanted in the thesis, extract all 24 \texttt{by\_mhz} entries
from \texttt{datasets/data\_multifreq\_train\_norm\_unbounded/normalization\_stats.json}}

\subsection{Impedance and frequency normalization}
\label{subsec:appA_imp_norm}

Impedance is normalized by a natural-log $z$-score,
$z = (\log(\max(Z, \texttt{1e-10})) - \texttt{log\_mean})/\texttt{log\_std}$, and saved per layout
as \texttt{imp.npy} of shape \texttt{(1, n\_freq)} in the normalized store.

\begin{table}[htbp]
\centering
\caption{\texttt{Impedance} and \texttt{PI\_freq} statistics blocks of the normalized training set.}
\label{tab:appA_imp_freq_stats}
\begin{tabular}{p{0.36\textwidth} p{0.50\textwidth}}
\hline
Key & Value \\
\hline
\texttt{Impedance.layout\_count} & 24979 \\
\texttt{Impedance.log\_mean} & \texttt{-1.039238452911377} \\
\texttt{Impedance.log\_std} & \texttt{1.7831873893737793} \\
\texttt{Impedance.z\_min} & \texttt{-2.7170498371124268} \\
\texttt{Impedance.z\_max} & \texttt{2.474900722503662} \\
\texttt{PI\_freq.min\_hz} & \texttt{1000000.0} \\
\texttt{PI\_freq.max\_hz} & \texttt{600000000.0} \\
\texttt{PI\_freq.log10\_min} & \texttt{6.0} \\
\texttt{PI\_freq.log10\_range} & \texttt{2.778151250383644} \\
\texttt{PI\_freq.count} & \texttt{688477} \\
\texttt{PI\_freq.anchor\_mhz} & length 24 \\
\texttt{PI\_freq.unique\_mhz\_rounded} & length 24 \\
\texttt{PI\_freq.expected\_anchors\_mhz} & length 24 \\
\texttt{multifreq\_manifest.rows} & 582997 \\
\texttt{multifreq\_manifest.unique\_design\_id} & 24979 \\
\texttt{multifreq\_manifest.max\_k\_filter} & 30 \\
\texttt{multifreq\_manifest} per-MHz counts & 24979 at every anchor except 80.0\,MHz with 8480 \\
\hline
\end{tabular}
\end{table}

\texttt{PI\_freq.count} (688477) equals the \emph{raw} manifest row count, while
\texttt{multifreq\_manifest.rows} in the same file is 582997, the normalized manifest row count.
This is recorded as an inconsistency in Section~\ref{sec:appA_inconsistencies} rather than resolved
here.

\section{Multi-type decap campaign (exp060) --- in flux}
\label{sec:appA_multitype}

The multi-type occupancy fork replaces the binary slot state with a per-slot type code. Its
simulation campaign was still running when the evidence for this appendix was first collected and
has since completed all 24 bands; the figures below were re-read on 2026-08-06 and are marked with
their source. What has \emph{not} happened is ingestion: no multi-type dataset directory was found
under \texttt{datasets/}, so nothing in this section is a dataset specification, and exp060 still
has no training input of its own.

\subsection{Type catalog and combination list}
\label{subsec:appA_multitype_catalog}

\begin{table}[htbp]
\centering
\caption{Multi-type decap catalog
(\texttt{scrap/generation/generate\_peb\_multitype.py}, \texttt{TYPE\_CATALOG}).}
\label{tab:appA_type_catalog}
\begin{tabular}{l r r r}
\hline
Type & $C$ & ESR & ESL \\
\hline
1 & \texttt{100e-9}\,F & \texttt{8.9e-3}\,$\Omega$ & \texttt{222e-12}\,H \\
2 & \texttt{47e-9}\,F & \texttt{21.4e-3}\,$\Omega$ & \texttt{154e-12}\,H \\
empty & \texttt{false} and \texttt{0} & --- & --- \\
\hline
\end{tabular}
\end{table}

The combination list \texttt{data/heatmaps/combinations\_multitype\_combined\_15000.csv} holds
15000 rows of 52 integer type codes (0 empty, 1 type-1, 2 type-2). Its sampler report
(\texttt{data/heatmaps/combinations\_multitype\_sample\_report.json}, seed 42, generated
\texttt{2026-08-02T18:07:25.393879+00:00}) records \texttt{per\_k}
$\{1{:}104,\ 2{:}5304,\ 3{:}4797,\ 4{:}3197,\ 5{:}1598\}$, families \texttt{type1\_only} 4576 /
\texttt{type2\_only} 4576 / \texttt{mixed} 5848, and \texttt{nonzero\_type\_counts}
$\{1{:}23015,\ 2{:}22866\}$. \texttt{docs/multitype\_peb\_pipeline.md} contradicts itself on these
numbers; see Section~\ref{sec:appA_inconsistencies}.

\subsection{Simulation pipeline configuration}
\label{subsec:appA_multitype_pipeline}

\texttt{pipelines/dataset\_sim/run\_multitype\_sim\_pipeline.py} reads the 15000-row CSV from row 0
with \texttt{MAX\_LAYOUTS} \texttt{None}, \texttt{USE\_ANCHOR\_MHZ} \texttt{False} and an explicit
\texttt{SWEEP\_MHZ} list of the same 24 anchors (10\ldots500\,MHz). \texttt{TRIGGER\_APPEND\_AFTER\_MOVE}
is \texttt{True}.

The append worker started by that pipeline
(\texttt{pipelines/dataset\_sim/run\_append\_worker.py}) runs
\texttt{pipelines/data/append\_merged\_combinations\_multifreq.py} with \texttt{OUTPUT\_ROOT} set to
\texttt{datasets/data\_multifreq\_train}, i.e.\ the single-type dataset root; that is the worker's
only output root. The worker log for the most recent jobs reports ``Tasks to run: 0'' and ``Nothing
to append''. This is stated as observed configuration, with no interpretation of intent.
\fillin{resolve the conflict between \texttt{TRIGGER\_APPEND\_AFTER\_MOVE=True} in
\texttt{run\_multitype\_sim\_pipeline.py} line 120 and the single-type-only
\texttt{OUTPUT\_ROOT} in \texttt{run\_append\_worker.py} lines 44--46; decide where multi-type raw
output is meant to be ingested and record the intended dataset root}

Append jobs are queued per MHz in \texttt{logs/append\_queue/pending.jsonl}, and completed jobs are
recorded in \texttt{logs/append\_queue/done.jsonl} with fields \texttt{mhz}, \texttt{append\_tag},
\texttt{raw\_root}, \texttt{enqueued\_at}, \texttt{started\_at}, \texttt{finished\_at} and
\texttt{returncode}. As read on 2026-08-06 the file holds 24 records, one per band, all with
\texttt{raw\_root}
\texttt{C:\textbackslash{}Users\textbackslash{}muthusamy\textbackslash{}Desktop\textbackslash{}Raw}.
The \texttt{merged\_470} record is dated 2026-07-05; the other 23 fall between 2026-08-03 and
2026-08-05. The last record is \texttt{mhz} 500, \texttt{append\_tag} \texttt{merged\_500},
\texttt{enqueued\_at} \texttt{2026-08-05T19:27:53.160450+00:00}, \texttt{finished\_at}
\texttt{2026-08-05T19:29:02.065218+00:00}, \texttt{returncode} 0. Twenty-three of the 24 records
carry \texttt{returncode} 0; \texttt{merged\_80} (\texttt{finished\_at}
\texttt{2026-08-03T17:38:40.736791+00:00}) carries \texttt{returncode} 1.

Three timestamps sit close together and are recorded here without an inference about which process
wrote what: \texttt{merged\_500} was enqueued at \texttt{2026-08-05T19:27:53.160450+00:00}, the
multi-type progress file was last updated at \texttt{2026-08-05T19:27:53.164372+00:00}
(\autoref{tab:appA_multitype_state}), and
\texttt{datasets/data\_multifreq\_train/dataset\_meta.json} carries \texttt{generated\_at\_utc}
\texttt{2026-08-05T19:28:42.489387+00:00} (\autoref{tab:appA_raw}).
\fillin{establish whether any multi-type raw output was appended into
\texttt{datasets/data\_multifreq\_train} by the append jobs of 2026-08-03 to 2026-08-05. The counts
in \autoref{tab:appA_raw} are unchanged from the earlier read, which does not by itself settle it;
the append worker's per-job log for \texttt{merged\_500}, and the \texttt{append\_tag} values in
\texttt{datasets/data\_multifreq\_train/manifest.csv}, would.}

\subsection{Campaign state}
\label{subsec:appA_multitype_state}

An earlier read of these artifacts caught the pipeline mid-run, with 23 of 24 bands complete. The
values below replace that snapshot: the simulation process
(\texttt{python -u pipelines/dataset\_sim/run\_multitype\_sim\_pipeline.py}) is no longer running,
and the progress file records all 24 bands as done. This is the state of the \emph{raw simulation
output on the Windows host}, not of a dataset.

\begin{table}[htbp]
\centering
\caption{Multi-type simulation state. The progress-file rows were re-read on 2026-08-06; the folder
counts were taken on 2026-08-05 after the final progress update.}
\label{tab:appA_multitype_state}
\begin{tabular}{p{0.40\textwidth} p{0.50\textwidth}}
\hline
Item & Value at the observation timestamp \\
\hline
Progress file & \texttt{C:\textbackslash{}Users\textbackslash{}muthusamy\textbackslash{}Desktop\textbackslash{}Raw\textbackslash{}sim\_progress\_multitype.json} \\
\texttt{impedance\_done} & \texttt{true} \\
\texttt{distribution\_mhz\_done} & 24 bands: 10, 63, 80, 100, 120, 150, 170, 180, 200, 230, 250,
270, 280, 300, 330, 350, 370, 390, 400, 420, 430, 450, 470, 500 --- the same 24 anchors as the
single-type set \\
\texttt{updated\_at} & \texttt{2026-08-05T19:27:53.164372+00:00} \\
Band folders inspected & the \texttt{Impedance} folder and the \texttt{heatmaps\_450MHz} and
\texttt{heatmaps\_500MHz} folders under the Raw root each contain 15000 PI folders, matching the
15000-row combination list \\
Multi-type dataset directory & none found under \texttt{datasets/} \\
\hline
\end{tabular}
\end{table}

\fillin{confirm the PI-folder count for the remaining 21 band folders under the Raw root; only
\texttt{Impedance}, \texttt{heatmaps\_450MHz} and \texttt{heatmaps\_500MHz} were counted, and a
completed band in \texttt{distribution\_mhz\_done} is not by itself proof of 15000 output folders}

\fillin{record the dataset root that the multi-type raw output is ingested into, once ingestion has
happened, together with the resulting counts, anchor list and normalization statistics; until then
exp060 has no dataset specification and no result from it can be reported}

\subsection{exp060 configuration}
\label{subsec:appA_exp060_config}

\texttt{experiments/exp060\_multitype\_occ/config.yaml} sets \texttt{n\_decap\_types} to 2 (line
193) while still pointing \texttt{data\_dir} at
\texttt{datasets/data\_multifreq\_train\_norm\_unbounded} (line 134), which is the single-type
normalized dataset. No multi-type normalized dataset directory was found under \texttt{datasets/}
at the read times given in \autoref{tab:appA_multitype_state}, so the completion of the simulation
campaign has not yet produced a training input for exp060.
\fillin{update or explain the exp060 \texttt{data\_dir}; it cannot supply the $52\times T$
multi-type occupancy that \texttt{n\_decap\_types = 2} implies}

\section{Recorded inconsistencies between artifacts}
\label{sec:appA_inconsistencies}

Five disagreements between repository artifacts were found while assembling this appendix. They are
listed rather than silently resolved, because in each case the authoritative source is either
ambiguous or the artifacts were written at different times.

\begin{table}[htbp]
\centering
\caption{Conflicts between artifacts. Both sides are named; none is resolved here.}
\label{tab:appA_conflicts}
\begin{tabular}{p{0.20\textwidth} p{0.34\textwidth} p{0.34\textwidth}}
\hline
Subject & Source A & Source B \\
\hline
Anchor count &
\texttt{docs/dataset.md}:96--101, ``Full 16-MHz set''; \texttt{configs/multifreq\_anchors.yaml}, 20
anchors &
\texttt{datasets/data\_multifreq\_train/dataset\_meta.json}, 24 anchors (the list actually loaded,
per the resolution order in \texttt{docs/dataset.md}:42--51) \\
\texttt{PI\_freq.count} &
\texttt{PI\_freq.count} = 688477, the raw manifest row count &
\texttt{multifreq\_manifest.rows} = 582997 in the same file,
\texttt{datasets/data\_multifreq\_train\_norm\_unbounded/normalization\_stats.json} \\
Artifact staleness &
normalized meta generated \texttt{2026-07-06T13:39:03}; 24979 layouts / 582997 rows &
raw meta generated \texttt{2026-08-05T19:28:42}; 29499 layouts / 688477 rows. The $K\le30$ filter
explains part of the gap, but the two are not from the same moment \\
Multi-type sampler &
\texttt{docs/multitype\_peb\_pipeline.md}:75--77: per-$K$ $\{2{:}3750, 3{:}3750, 4{:}3750,
5{:}3750\}$, families 5000 each, nonzero type counts $\{1{:}26258, 2{:}26242\}$ &
\texttt{data/heatmaps/combinations\_multitype\_sample\_report.json} and the CSV on disk: per-$K$
$\{1{:}104, 2{:}5304, 3{:}4797, 4{:}3197, 5{:}1598\}$, families 4576/4576/5848, nonzero
$\{1{:}23015, 2{:}22866\}$. The same doc at lines 66--68 quotes the report's numbers, so it
contradicts itself \\
Merged CSV path &
\texttt{merged\_combinations\_report.json} key \texttt{output\_csv} and
\texttt{docs/dataset.md}:78--81 point to
\texttt{data/heatmaps/all\_combinations\_merged.csv}, which is not present there &
the 29499-row merged CSV on disk is
\texttt{data/heatmaps/peb\_with\_29k/all\_combinations\_merged.csv}, which is what
\texttt{run\_combinations\_sim\_pipeline.py}:50 reads \\
Append worker root &
\texttt{run\_multitype\_sim\_pipeline.py}:120 sets
\texttt{TRIGGER\_APPEND\_AFTER\_MOVE=True} for the multi-type campaign &
\texttt{run\_append\_worker.py}:44--46 has \texttt{OUTPUT\_ROOT =
REPO\_ROOT/datasets/data\_multifreq\_train}, the single-type dataset, as its only output root \\
\hline
\end{tabular}
\end{table}

\fillin{decide, for each row of Table~\ref{tab:appA_conflicts}, which artifact is authoritative and
either correct the stale one in the repository or state the resolution in the thesis text}

\section{Reproducibility notes}
\label{sec:appA_repro}

Two constraints limit reproduction of the work in this appendix. First, the simulation half is
bound to one Windows host: \texttt{engineer.exe}, the \texttt{.erf} design, the PEB staging
directories and the Raw output root are all absolute
\texttt{C:\textbackslash{}Users\textbackslash{}muthusamy\textbackslash{}\ldots} paths, and the AL
loop cannot run its simulation steps without that machine. Second, every dataset directory is
git-ignored, so the repository alone does not carry the data.

\fillin{record a wall-clock simulation cost per layout and per band; \texttt{docs/ecadstar\_headless\_cli.md}
lines 113--121 gives an observed rate for one resumed band, but no systematic timing artifact
exists. A per-band timing table from \texttt{logs/append\_queue/done.jsonl} start/finish timestamps
would settle it}

\fillin{state whether \texttt{H-shape.erf} and the derived datasets can be released, and under what
terms, so the reproducibility statement is complete}

% ---------------------------------------------------------------
% FILL IN checklist:
%   - Port convention (IC1_Port1 vs IC1_Port1,IC2_Port2) behind the stored 231-bin spectra:
%     check harvested CSV names in a legacy PI folder vs peb_components of the producing pipeline.
%   - Solver settings block (Table tab:appA_solver_gaps): sweep definition, sweep points,
%     mesh/plane-pair settings, port definitions, stack-up, materials, board dimensions,
%     single-type decap C/ESR/ESL. All live in H-shape.erf on the Windows host; export them.
%   - Whether 490/550/590 MHz from SWEEP_MHZ ever reached the dataset (logs/append_queue/done.jsonl).
%   - Provenance of configs/Frequency_data_hz.npy (how the 231-point grid was generated).
%   - Provenance of configs/target_impedance.npy (per-band values, derivation, 0.8 / 50 limits).
%   - Train/validation/test split definition for the normalized dataset (exp059 dataloader).
%   - Reason the 80 MHz band covers 10000 (raw) / 8480 (normalized) layouts instead of the pool.
%   - Provenance of the 24-anchor selection.
%   - Anchor-count conflict: docs/dataset.md "16-MHz set" vs 20-anchor YAML vs 24-anchor meta.
%   - Full k_hist_all histogram, if the budget distribution is discussed in the main text.
%   - Full 24-entry per-MHz heatmap statistics table, if wanted.
%   - Append-root conflict for the multi-type campaign (TRIGGER_APPEND_AFTER_MOVE vs OUTPUT_ROOT).
%   - Whether multi-type raw output was appended into datasets/data_multifreq_train by the
%     2026-08-03..05 append jobs (per-job worker log for merged_500; append_tag values in manifest.csv).
%   - PI-folder counts for the 21 band folders that were not counted (only Impedance, 450 and 500 MHz were).
%   - Dataset root that the multi-type raw output is ingested into, plus its counts/anchors/stats.
%   - Whether SWEEP_MHZ in run_combinations_sim_pipeline.py was edited between campaigns.
%   - exp060 config data_dir vs n_decap_types = 2 (no multi-type dataset exists).
%   - Per-row resolution of Table tab:appA_conflicts.
%   - Wall-clock simulation cost per layout and per band.
%   - Release/availability terms for H-shape.erf and the datasets.
%
% Citation verification queue:
%   - None. This appendix cites no literature by instruction; every statement is sourced from a
%     repository artifact named in the surrounding text.
%
% Evidence warnings:
%   - Section sec:appA_multitype reports the Windows simulation host, not repository artifacts.
%     The campaign has FINISHED: sim_progress_multitype.json lists all 24 bands with updated_at
%     2026-08-05T19:27:53.164372+00:00 (re-read 2026-08-06) and the pipeline process is gone.
%     Folder counts (15000) were taken on 2026-08-05 for Impedance, 450 MHz and 500 MHz only.
%     No multi-type dataset directory was found under datasets/; exp060 still has no dataset
%     specification and no exp060 result may be reported.
%   - datasets/data_multifreq_train is written to by the append worker. generated_at_utc moved
%     from 2026-08-05T15:24:43.778950+00:00 to 2026-08-05T19:28:42.489387+00:00 between reads
%     while all counts stayed identical. tab:appA_raw carries the later value; re-read before
%     submission. Whether multi-type output entered that root is an open question, flagged in
%     subsec:appA_multitype_pipeline.
%   - logs/append_queue/done.jsonl: merged_80 (2026-08-03T17:38:40.736791+00:00) has returncode 1.
%     Not linked here to the 80 MHz coverage shortfall in tab:appA_raw - that link is unverified.
%   - Five artifact conflicts are recorded in Table tab:appA_conflicts and none is resolved.
%     Do not quote one side of any of them in the main text without the other.
%   - datasets/data_multifreq_train_norm_unbounded/dataset_meta.json (2026-07-06) predates the raw
%     set it derives from (2026-08-05). Any count taken from it is a snapshot, not current state.
%   - datasets/data_multifreq_al_overlay_exp059 has no dataset_meta.json; its counts were read from
%     a directory listing, so they carry no generation timestamp.
%   - configs/multifreq_anchors.yaml is a legacy fallback (20 anchors) and must not be used as the
%     anchor list in any chapter.
%   - Anchor resolution order is quoted from docs/dataset.md; the implementation in
%     src_vae/others/multifreq_anchors.py was not read line by line.
%   - SWEEP_MHZ in run_combinations_sim_pipeline.py is a plain module constant with no history in
%     the repository; no claim is made about when it was last edited or which sweep it reflects.
%   - dataset_root and source_dataset in dataset_meta.json are absolute paths (libs/dataset_meta.py
%     line 172), not repository-relative ones.
%   - No result, accuracy, or performance claim appears in this appendix by design.
% ---------------------------------------------------------------
